\section{Introduction}\label{sec:introduction}

Birational anabelian geometry asks how much of a function field can be
recovered from its Galois-theoretic invariants. In the Bogomolov program,
the field has algebraically closed constants, and the relevant invariants
are much smaller than the full absolute Galois group. The present paper
proves the reconstruction assertion for the first two Milnor $K$-groups
modulo a single prime, with their multiplication. The result includes
the isomorphism statement: it identifies every compatible linear isomorphism,
as well as the ambiguity in the field isomorphism that induces it.

\subsection{The reconstruction theorem}

Fix a prime $\ell$ and put $\Lambda=\F_\ell$. For a field $F$ of
characteristic different from $\ell$, set
\[
 V_F=F^\times/(F^\times)^\ell,
 \qquad
 R_F=\Span_\Lambda\{[x]\otimes[1-x]:x\in F\setminus\{0,1\}\}.
\]
We write $V_F$ additively, so $[xy]=[x]+[y]$. Let
\[
 W_F=(V_F\otimes_\Lambda V_F)/R_F,
 \qquad m_F:V_F\otimes_\Lambda V_F\longrightarrow W_F.
\]
Thus $V_F=K^{\mathrm M}_1(F)/\ell$ and
$W_F=K^{\mathrm M}_2(F)/\ell$, with $m_F$ the Milnor product.
A $\Lambda$-linear isomorphism $\Theta:V_K\to V_L$ is
\emph{compatible} if
\[
 (\Theta\otimes\Theta)(R_K)=R_L.
\]
Equivalently, there is a unique linear isomorphism $\Theta_2:W_K\to W_L$
intertwining the products. Write $\Isom_{\mathrm M}(V_K,V_L)$ for
the set of compatible isomorphisms. Multiplication by
$a\in\Lambda^\times$ acts on this set; the induced degree-two map
is then multiplied by $a^2$.

Let $K/k$ and $L/l$ be finitely generated extensions of algebraically
closed fields. Write $F^i$ for the perfect closure of a field $F$:
it equals $F$ in characteristic zero and
$\bigcup_{r\ge0}F^{1/p^r}$ in characteristic $p>0$. Denote by
$\Isom^i(K,L)$ the field isomorphisms $\alpha:K^i\to L^i$
satisfying $\alpha(k)=l$. In positive characteristic, identify two
such isomorphisms if they differ by postcomposition with
$\Frob_L^n$, $n\in\Z$, where $\Frob_L(z)=z^p$. Denote the resulting
set by $\Isom^i_{\mathrm F}(K,L)$; in characteristic zero take no
quotient, and if the characteristics differ this set is empty.

Purely inseparable extension induces canonical isomorphisms on $V$
and $W$. Indeed, for an extension of exponent at most $r$, inclusion
and the $p^r$-power homomorphism compose to the $p^r$-power map on
each field. On Milnor groups in degrees one and two this composition
is multiplication by $p^r$ and $p^{2r}$, respectively, both invertible
modulo $\ell$. Passing to the union proves the assertion for perfect
closures. Thus $\alpha$ induces a compatible map $\alpha_1:V_K\to V_L$,
and postcomposition by $\Frob_L^n$ multiplies it by $p^n$ in
$\Lambda^\times$. There is consequently a canonical map
\begin{equation}\label{eq:main-map}
 \Isom^i_{\mathrm F}(K,L)\longrightarrow
 \Isom_{\mathrm M}(V_K,V_L)/\Lambda^\times.
\end{equation}

\begin{theorem}\label{thm:main}
Let $\ell$ be any prime. Let $K/k$ and $L/l$ be finitely generated
extensions of arbitrary algebraically closed fields, with
$\charac k,\charac l\ne\ell$ and
$\trdeg(K/k),\trdeg(L/l)\ge2$. Then \eqref{eq:main-map} is bijective.
In particular, a compatible isomorphism $V_K\to V_L$ forces equality
of the characteristics and of the relative transcendence degrees.
\end{theorem}

Only the two vector spaces and their bilinear product enter this
statement. The proof recovers the valuations and rational-subfield
images needed for reconstruction from this datum. The scalar is one
global element of $\F_\ell^\times$; at $\ell=2$ this ambiguity is trivial.

\subsection{Historical context}

Bogomolov proposed reconstructing higher-dimensional function fields over
algebraically closed constants from the pro-$\ell$ quotient in which
commutators are central \cite{Bogomolov1991}. Bogomolov and Tschinkel
proved reconstruction for surfaces over algebraic closures of finite
fields \cite{BT2008}, and subsequently for higher-dimensional function
fields over those constants \cite{BT2011}. Pop independently completed
reconstruction in this constant-field setting \cite{PopBI2012}.
These results established that a small part of the Galois group can
retain the geometry needed to recover a field. Over more general
algebraically closed constants, Pop also obtained reconstruction from
pro-$\ell$ data endowed with divisorial inertia, in relative
transcendence degree greater than two \cite[Theorem~1.1]{Pop2018}.

The finite-coefficient problem requires additional control of the global
reconstruction. Topaz proved a Milnor-theoretic isomorphism theorem in
relative transcendence degree at least five when the degree-one and
degree-two mod-$\ell$ groups and their product are supplied together
with all rational subgroups \cite[Theorem~B]{Topaz2016}. Those subgroups
are the images of the multiplicative groups of relatively algebraically
closed rational one-variable subfields. Theorem~\ref{thm:main} obtains
the required rational subfields from the product itself, and its
point-value argument applies in every relative dimension at least two.

Algebraic dependence has provided another route from multiplicative
invariants to field structure. Bogomolov and Tschinkel used Milnor
$K$-theory modulo infinitely divisible elements in characteristic zero
\cite{BT2009}. Cadoret and Pirutka reconstructed regular function
fields over perfect constants from the multiplicative quotient by
constants together with algebraic dependence, and derived applications
to integral Milnor $K$-theory \cite{CadoretPirutka2021}. Topaz later
proved reconstruction from rational Milnor $K$-theory in absolute
transcendence degree at least five \cite{Topaz2022}. These results
clarify the role of algebraic dependence, while using invariants
different from the single-prime datum considered here.

Our local inputs are the alternating-pair valuation theory recorded in
\cite[Section~2]{Topaz2020} and Pop's density theorem for minimized
inertia \cite{Pop2015}. We state the precise forms used below.
The bounded-support and incidence method is adapted from
\cite[Sections~4--6]{InternalBP}. The finite-coefficient bounded-support
and incidence arguments are proved here. The support estimate also has an
antecedent in the bounded-genus curve sections and Hurwitz argument of
Bogomolov and Tschinkel \cite[Section~6, Proposition~6.1]{BT2011}.
The proof uses no pro-$\ell$ field-reconstruction theorem as a premise.

\subsection{The proof and its main ingredients}

The argument begins with a compatible map $\Theta$ and its dual on
the compact character spaces
$A_F=\Hom(F^\times,\Lambda)$. Two characters $f,g$ form an
\emph{alternating pair} when
$f(x)g(1-x)=f(1-x)g(x)$ for every $x\ne0,1$. This condition is
visible in the Milnor product and is therefore preserved by $\Theta$.
Established local theory turns alternating subspaces into valuations.
Section~\ref{sec:local} uses it to recover the transcendence degree and
the inertia and decomposition spaces of quasi-divisorial valuations,
and the character spaces of their iterated residue fields. These
valuations may act nontrivially on the constants; their residue
constants can therefore have positive characteristic even when the
original field has characteristic zero.

The next task is to recover curve subfields: relatively algebraically
closed intermediate fields $k\subset E\subset K$ of transcendence
degree one over $k$. Individual
mod-$\ell$ classes record only orders modulo $\ell$, so a divisor can
disappear when its multiplicity is divisible by $\ell$. We control
this loss by passing to residual curves. Section~\ref{sec:residual-curves}
uses Pop's density theorem to recognize residual curves whose constants
are algebraic closures of finite fields, together with all their point
order lines. Section~\ref{sec:support} then constructs a residual
curve of bounded genus for any finite independent list of classes
in $V_L$. It preserves independence and realizes each support degree
on a fixed projective model of $L$ as the number of nonzero point
orders on the residual curve.
For each nonconstant $t\in K$, applying Riemann--Hurwitz to five fixed
members of its pencil bounds the total degree of the divisor support
of every class $\Theta([t-a])$, $a\in k$,
on that model.

For a curve subfield $E$, choose $t\in E\setminus k$. The supporting
prime divisors of the classes $\Theta([t-a])$ form an infinite set of
bounded degree, hence lie in a finite-type family. The valuation
correspondence shows that each class in
$\Theta(V_E)$ has an $\ell$-th-power residue along all but finitely many
of these divisors. Section~\ref{sec:curve-fields} studies the incidence
variety of pairs consisting of a point and a divisor containing it.
Restricting its parameter space to a curve produces a finite extension
of $L$ containing the curve's function field. After passage to this
extension, the classes come from that one-variable parameter field.
The argument then descends them to a curve subfield $P\subset L$
and proves $\Theta(V_E)=V_P$.

Two features are useful specifically for finite coefficients. One
parameter curve works for every class. For each class, a finite change
of parameter space allows normality to extend a generic $\ell$-th root
of a representative across the smooth locus where that representative
is a unit. This
avoids any restriction on the cardinality of the constants. Moreover,
the incidence field is the compositum of $L$ and the parameter field.
This identity permits exact descent, retaining information that a
norm could annihilate modulo $\ell$.

Once curve subfields correspond, Section~\ref{sec:good-functions}
recognizes true divisors, which are trivial on the constants, and a
sufficiently large family of rational
subfields on which every point is detected. A quotient of two local
parameters at a smooth point supplies such a subfield. These quotients
generate the function field, and the construction works already in
dimension two. Matching their point-order lines gives bijections
between the corresponding projective lines of points.

Section~\ref{sec:field} recovers field operations from simultaneous
values of these rational functions. A blowup at a smooth point detects
all coordinates of a finite tuple at once. Additive triples yield
rational expressions, allowing $p$-power roots in characteristic $p$,
for the transported addition law.
Their one-variable slices align all point bijections, outside finite
sets, with a single isomorphism of constant fields. An identity of rational functions
removes the apparent freedom caused by finite exceptional sets in
those slices. The resulting correspondence preserves every algebraic
relation among the chosen generators and hence gives an isomorphism
of perfect fields.

Finally, Section~\ref{sec:compatibility} proves that this isomorphism
induces the original map up to one scalar, and that scalar action of
a field automorphism forces it to be a Frobenius power. Throughout,
the pro-$\ell$ density input concerns individual fields. The given
mod-$\ell$ isomorphism is never lifted to a pro-$\ell$ isomorphism.

\subsection{Conventions}

A function field $F/\kappa$ means a finitely generated field extension.
Its constants $\kappa$ will be algebraically closed. A curve is a
geometrically integral variety of dimension one, and its function field
has a unique smooth projective model over algebraically closed constants.
A field extension is regular if it is separable and its base is
relatively algebraically closed. Points of varieties over an
algebraically closed field mean rational points unless a scheme point
is specified. The phrase ``for general points'' means on some dense
Zariski open subset. We continue to use
$V_F=F^\times/(F^\times)^\ell$ for residue fields of characteristic
$\ell$, as a multiplicative quotient, without a Galois-cohomological
interpretation.
